\section{Визуализация светящейся среды}

При синтезе в узлах расчётной сетки вычисляется плотность светящейся среды $\rho \in [0;\, 1]$: коэффициент излучения в точке равен $\varepsilon \rho$, коэффициент поглощения "--- $\kappa \rho$, а цвет $C(z)$ определяется высотой. По закону Бугера интенсивность света, прошедшего в поглощающей среде путь $l$, равна~\cite{landsberg2003}
\begin{equation}
	I = I_0\, e^{-\kappa l}.
	\label{eq:bouguer}
\end{equation}
Таким образом, на пути $ds$ интенсивность убывает на долю $\kappa \rho\, ds$, а среда сама излучает $\varepsilon \rho C\, ds$, поэтому яркость, приходящая к наблюдателю вдоль луча зрения, равна
\begin{equation}
	I = \int_0^{s_{\max}} \varepsilon \rho(s)\, C(z(s))\, T(s)\, ds + I_{\text{ф}}\, T(s_{\max}),
	\label{eq:volume}
\end{equation}
где
\begin{equation}
	T(s) = \exp\left( -\int_0^{s} \kappa \rho(s')\, ds' \right)
	\label{eq:transmittance}
\end{equation}
"--- пропускание среды, $s_{\max}$ "--- расстояние до ближайшей непрозрачной поверхности или до выхода из области сетки, $I_{\text{ф}}$ "--- яркость фона.

\subsection{Полупрозрачные срезы}

Область сетки изображается параллельными полупрозрачными плоскостями на расстоянии $\Delta s$, которые смешиваются от дальней к ближней по правилу $I = k I_1 + (1 - k) I_2$, где $k$ "--- непрозрачность ближней плоскости~\cite[с.~413]{rogers1989}. Это приближённое вычисление интеграла~\eqref{eq:volume}, но при угле $\gamma$ между лучом и нормалью к срезам шаг вдоль луча возрастает до $\Delta s / \cos\gamma$.

\subsection{Трассировка лучей в среде}

Через каждый пиксель из точки наблюдения строится луч~\cite[с.~360--361]{rogers1989}
\begin{equation}
	P(s) = P_V + s \vec{D}, \qquad s \geqslant 0,
	\label{eq:ray}
\end{equation}
где $\vec{D}$ "--- единичный вектор направления через пиксель. Вдоль луча интеграл~\eqref{eq:volume} вычисляется с шагом $\Delta s$; значение $\rho$ получается трилинейной интерполяцией восьми соседних узлов сетки:
\begin{equation}
	I \leftarrow I + T \varepsilon \rho\, C\, \Delta s, \qquad T \leftarrow T \exp(-\kappa \rho\, \Delta s).
	\label{eq:march}
\end{equation}
Луч обрывается на непрозрачной поверхности (её глубина известна из z-буфера), на выходе из сетки или при пренебрежимо малом $T$.

\subsection{Выбор метода визуализации светящейся среды}

Число выборок на кадр для обоих методов $\approx p \cdot L_{\text{л}} / \Delta s$, где $L_{\text{л}}$ "--- длина пути луча в сетке; при $L_{\text{л}} = 100$~у.е. и $\Delta s = 0{,}4$~у.е. это около $5 \cdot 10^8$. Сравнение приведено в таблице~\ref{tab:volume}.

\begin{table}[H]
	\caption{Сравнение методов визуализации светящейся среды}
	\label{tab:volume}
	\begin{tabular}{|L{0.44}|C{0.27}|C{0.27}|}
		\hline
		Критерий & Полупрозрачные срезы & Трассировка лучей \\
		\hline
		Шаг вдоль луча & $\Delta s / \cos\gamma$ & $\Delta s$ \\
		\hline
		Выборок на кадр & $\approx 5 \cdot 10^8$ & не более $5 \cdot 10^8$ \\
		\hline
		Досрочное завершение луча & нет & да \\
		\hline
	\end{tabular}
\end{table}

У срезов шаг вдоль луча зависит от направления взгляда, поэтому выбрана трассировка лучей в среде~\eqref{eq:march}.